% BHU PYQ -- Manually transcribed & error-corrected
% Paper: BCF-321 : Portfolio Management
% Exam: B.Com. (Hons.) FMM Semester VI Examination, 2022-23

\documentclass[11pt,a4paper]{article}
\usepackage[a4paper,top=2.5cm,bottom=2.5cm,left=2.8cm,right=2.8cm,headheight=14pt]{geometry}
\usepackage{amsmath,amssymb}
\usepackage[T1]{fontenc}
\usepackage[utf8]{inputenc}
\usepackage{lmodern,microtype,fancyhdr,enumitem,hyperref,lastpage,parskip}

\hypersetup{
    colorlinks=true,
    urlcolor=black,
    linkcolor=black,
    pdftitle={BCF-321: Portfolio Management -- BHU B.Com. (Hons.) FMM Sem VI 2022-23}
}

\pagestyle{fancy}
\fancyhf{}
\renewcommand{\headrulewidth}{0.4pt}
\renewcommand{\footrulewidth}{0.4pt}
\fancyhead[L]{\small\textit{Banaras Hindu University}}
\fancyhead[R]{\small\textit{BCF-321: Portfolio Management}}
\fancyfoot[C]{\small Page \thepage\ of \pageref{LastPage}}

\newlist{parts}{enumerate}{2}
\setlist[parts,1]{label=(\alph*),leftmargin=*,itemsep=5pt,topsep=3pt}
\setlist[parts,2]{label=(\roman*),leftmargin=*,itemsep=3pt,topsep=2pt}
\newcommand{\pts}[1]{\hfill\textit{[#1]}}

\begin{document}

\begin{center}
  {\large\bfseries BANARAS HINDU UNIVERSITY}\\[2pt]
  {\normalsize B.Com. (Hons.) FMM --- Semester VI Examination, 2022-23}\\[6pt]
  \rule{0.8\linewidth}{0.5pt}\\[6pt]
  {\large\bfseries Paper No. BCF-321: Portfolio Management}\\[4pt]
  {\normalsize Time: 3 Hours\hfill Maximum Marks: 70}
\end{center}

\rule{\linewidth}{0.4pt}

\smallskip
\noindent\textit{Instructions: Attempt all questions. All questions carry equal marks (14 marks each).}

\bigskip

\noindent\textbf{Question 1.}
Present your understanding of ``Portfolio of Securities''. Also explain different phases of Portfolio Management.\pts{14}

\centerline{\textbf{OR}}

Write a note on evolution of Portfolio Management. In what ways it has become more relevant nowadays?\pts{14}

\medskip\rule{\linewidth}{0.2pt}

\noindent\textbf{Question 2.}
Elucidate the component of ``Risk'' in detail in reference to Portfolio Management.\pts{14}

\centerline{\textbf{OR}}

Past data for four years on two securities `A' and `B' show that the securities have yielded return on investment as under:
\begin{center}
\begin{tabular}{|c|c|c|}
\hline
\textbf{Year} & \textbf{Return from Security A (\%)} & \textbf{Return from Security B (\%)} \\ \hline
1998--1999 & 15 & 9 \\ \hline
1999--2000 & 5 & 11 \\ \hline
2000--2001 & 18 & 15 \\ \hline
2001--2002 & 7 & 10 \\ \hline
\end{tabular}
\end{center}
Calculate the average return offered by the two securities and their relative riskiness.\pts{14}

\medskip\rule{\linewidth}{0.2pt}

\noindent\textbf{Question 3.}
A portfolio has four securities and the expected returns from them are:
\begin{align*}
  r_1 &= 15\% \\
  r_2 &= 12\% \\
  r_3 &= 14\% \\
  r_4 &= 20\%
\end{align*}
The funds invested in four securities are Rs. 2,00,000, Rs. 2,80,000, Rs. 3,20,000 and Rs. 4,00,000 respectively. Find the expected return from the portfolio.\pts{14}

\centerline{\textbf{OR}}

Compare and contrast ``Fundamental Analysis'' with ``Technical Analysis''.\pts{14}

\medskip\rule{\linewidth}{0.2pt}

\noindent\textbf{Question 4.}
Following information is available in respect of Market:
\begin{center}
\begin{tabular}{|c|c|c|}
\hline
\textbf{Security} & \textbf{Expected Return (\%)} & \textbf{Beta ($\beta$)} \\ \hline
1 & 22.20 & 1.75 \\ \hline
2 & 15.80 & 1.90 \\ \hline
3 & 18.00 & 1.10 \\ \hline
4 & 9.00 & 0.95 \\ \hline
5 & 25.80 & 2.00 \\ \hline
Risk-Free Asset & 8.00 & --- \\ \hline
Market Portfolio & 15.00 & 1.00 \\ \hline
\end{tabular}
\end{center}
Find out which of the securities are overpriced or underpriced.\pts{14}

\centerline{\textbf{OR}}

Explain any one Model used for ``Portfolio Selection'' in detail.\pts{14}

\medskip\rule{\linewidth}{0.2pt}

\noindent\textbf{Question 5.}
Compare the following two portfolios on the basis of Treynor ratio and Sharpe ratio:
\begin{center}
\begin{tabular}{|c|c|c|c|}
\hline
\textbf{Portfolio} & \textbf{Return (\%)} & \textbf{Standard Deviation (\%)} & \textbf{Beta ($\beta$)} \\ \hline
A & 10 & 13 & 0.75 \\ \hline
B & 20 & 26 & 1.45 \\ \hline
Market Portfolio & 14 & 18 & 1.00 \\ \hline
\end{tabular}
\end{center}
Interest-free rate of return = 8\%.\pts{14}

\centerline{\textbf{OR}}

Give a detailed insight of strategies and constraints of Portfolio Revision.\pts{14}

\vfill
\rule{\linewidth}{0.4pt}
\begin{center}\small
\href{https://bhu-exam.vercel.app/}{Click here to visit our website}\\[4pt]
\href{https://me-aryan.vercel.app/}{Click here to visit our contributor}\\[4pt]
\href{https://quashan-me.vercel.app/}{Click here to visit our contributor}
\end{center}

\end{document}
